\chapter{System Architecture}

\section{Architecture Overview}

D-CARPOOL follows a three-tier architecture comprising:

\begin{itemize}
    \item \textbf{Presentation Layer}: React-based frontend with responsive design and MetaMask integration
    \item \textbf{Application Layer}: Node.js/Express backend with RESTful APIs and cryptographic services
    \item \textbf{Data Layer}: MySQL database, Ethereum blockchain, and IPFS for decentralized storage
\end{itemize}

\begin{figure}[H]
\centering
\fbox{\parbox{0.9\textwidth}{
\centering
\vspace{0.5cm}
\textbf{Three-Tier Architecture}\\[0.3cm]
\begin{tabular}{|c|}
\hline
\textbf{PRESENTATION LAYER} \\
React.js + TailwindCSS + MetaMask + Leaflet \\
\hline
\end{tabular}\\[0.3cm]
$\downarrow$ REST API / WebSocket $\downarrow$\\[0.3cm]
\begin{tabular}{|c|}
\hline
\textbf{APPLICATION LAYER} \\
Node.js + Express + Sequelize + ethers.js \\
\hline
\end{tabular}\\[0.3cm]
$\downarrow$\hspace{2cm}$\downarrow$\hspace{2cm}$\downarrow$\\[0.3cm]
\begin{tabular}{|c|c|c|}
\hline
MySQL & Ethereum & IPFS \\
Database & Blockchain & Storage \\
\hline
\end{tabular}
\vspace{0.5cm}
}}
\caption{System Architecture of D-CARPOOL}
\end{figure}

\section{Technology Stack}

\subsection{Frontend Technologies}

\begin{itemize}
    \item \textbf{React.js 18}: Component-based UI framework with hooks
    \item \textbf{TailwindCSS 3}: Utility-first CSS framework for responsive design
    \item \textbf{Ethers.js 6}: Ethereum library for blockchain interaction
    \item \textbf{React Router 6}: Client-side routing for SPA navigation
    \item \textbf{Axios}: HTTP client for REST API requests
    \item \textbf{React-Leaflet 4.2.1}: OpenStreetMap integration for maps
    \item \textbf{React Hot Toast}: User notification system
    \item \textbf{Lucide React}: Modern icon library
\end{itemize}

\subsection{Backend Technologies}

\begin{itemize}
    \item \textbf{Node.js 18}: JavaScript runtime environment
    \item \textbf{Express.js 4}: Web application framework
    \item \textbf{MySQL 8}: Relational database management system
    \item \textbf{Sequelize 6}: Promise-based ORM for Node.js
    \item \textbf{Nodemailer}: Email service integration for OTP
    \item \textbf{Bcrypt}: Password hashing library
    \item \textbf{secrets.js-grempe}: Shamir's Secret Sharing library
    \item \textbf{merkletreejs}: Merkle tree implementation
\end{itemize}

\subsection{Blockchain Technologies}

\begin{itemize}
    \item \textbf{Solidity 0.8.20}: Smart contract programming language
    \item \textbf{Hardhat}: Ethereum development environment
    \item \textbf{OpenZeppelin}: Secure smart contract library
    \item \textbf{IPFS/Pinata}: Distributed file storage
\end{itemize}

\subsection{Cryptographic Services}

\begin{itemize}
    \item \textbf{CP-ABE Service (Flask)}: Attribute-based encryption (Port 7001)
    \item \textbf{Shamir SSS Utility}: Secret splitting and reconstruction
    \item \textbf{Merkle Tree Utility}: Batch anchoring and proof generation
\end{itemize}

\section{Database Schema}

The database consists of eight primary tables:

\begin{enumerate}
    \item \textbf{users}: Stores user account information and wallet addresses
    \item \textbf{email\_verifications}: Manages OTP verification
    \item \textbf{rides}: Contains ride details and agreement hashes
    \item \textbf{ride\_participants}: Tracks ride participation and status
    \item \textbf{ratings}: Stores user ratings and reviews
    \item \textbf{complaints}: Manages user complaints with voting
    \item \textbf{sos\_alerts}: Records emergency alerts with share distribution
    \item \textbf{tickets}: Manages PNR-based train ticket integration
\end{enumerate}

\subsection{Entity-Relationship Model}

The relationships between entities are defined as:

\begin{itemize}
    \item A User can create multiple Rides (one-to-many)
    \item A User can participate in multiple Rides (many-to-many through ride\_participants)
    \item A User can give multiple Ratings (one-to-many)
    \item A User can file multiple Complaints (one-to-many)
    \item A Ride can have multiple Participants (one-to-many)
    \item A Ride can have multiple Ratings (one-to-many)
    \item A User can have multiple SOS configurations (one-to-many)
\end{itemize}

\section{Smart Contract Architecture}

The blockchain layer consists of eight smart contracts:

\begin{table}[H]
\centering
\caption{Smart Contract Deployment}
\begin{tabular}{|l|l|l|}
\hline
\textbf{Contract} & \textbf{Purpose} & \textbf{Address} \\
\hline
UserIdentity & Soulbound Token Management & 0xc6e7...4e7d \\
RideContract & Ride Agreement Hashes & 0x59b6...857b \\
Reputation & On-chain Ratings & 0x4ed7...e1C1 \\
DPoSGovernance & Delegate Elections & 0x3228...6c44 \\
DisputeResolution & Dispute Voting & 0xa852...38f \\
ZKPDriverVerifier & ZK Proof Verification & 0x4A67...4319 \\
MerkleAuditTrail & Batch Anchoring & 0x7a20...814F \\
ShamirSOSVault & SSS Share Commitments & 0x0963...eBef \\
\hline
\end{tabular}
\end{table}

\subsection{UserIdentity Contract}

Manages Soulbound Token (SBT) for academic identity:

\begin{lstlisting}[language=Java, caption=UserIdentity Contract Structure]
struct User {
    address walletAddress;
    string username;
    string email;
    string ipfsHash;
    bool isActive;
    uint256 registeredAt;
    bool hasSBT;  // Soulbound Token status
}
\end{lstlisting}

Key functions:
\begin{itemize}
    \item \texttt{mintSBT()}: Mint non-transferable academic identity token
    \item \texttt{verifySBT()}: Check if address has valid SBT
    \item \texttt{revokeSBT()}: Admin revocation for policy violations
\end{itemize}

\subsection{RideContract}

Manages ride lifecycle and agreement hashes:

\begin{lstlisting}[language=Java, caption=Ride Agreement Structure]
struct Ride {
    uint256 rideId;
    address host;
    bytes32 agreementHash;  // keccak256 of ride details
    uint8 availableSeats;
    RideStatus status;
    uint256 createdAt;
}
\end{lstlisting}

\subsection{ShamirSOSVault Contract}

Manages emergency data protection with threshold sharing:

\begin{lstlisting}[language=Java, caption=SOS Vault Structure]
struct SOSConfig {
    address user;
    uint8 totalShares;      // n = 5
    uint8 threshold;        // k = 3
    bytes32[] shareHashes;  // Commitments
    bool isActive;
    uint256 configuredAt;
}
\end{lstlisting}

\section{API Architecture}

The backend exposes RESTful APIs organized into eight modules:

\begin{table}[H]
\centering
\caption{API Endpoint Categories}
\begin{tabular}{|l|c|l|}
\hline
\textbf{Module} & \textbf{Endpoints} & \textbf{Purpose} \\
\hline
Authentication & 4 & User registration, login, wallet linking \\
Email Verification & 4 & OTP generation, verification, resend \\
Rides & 10 & CRUD operations, search, join/leave \\
Ratings & 4 & Submit ratings, get user reputation \\
Complaints & 5 & File complaints, delegate voting \\
SOS & 6 & Trigger alerts, SSS operations \\
Admin & 8 & Platform administration, Merkle anchoring \\
IPFS & 3 & Document upload, retrieval, verification \\
\hline
\end{tabular}
\end{table}

\section{Security Architecture}

\subsection{Multi-Layer Security Model}

D-CARPOOL implements three security layers:

\begin{enumerate}
    \item \textbf{Layer 1 - Identity \& Access Control}:
    \begin{itemize}
        \item Soulbound Tokens for verified academic identity
        \item JWT authentication with access/refresh tokens
        \item Email OTP verification for @iiitkottayam.ac.in
        \item MetaMask wallet connection for blockchain operations
    \end{itemize}
    
    \item \textbf{Layer 2 - Data Protection}:
    \begin{itemize}
        \item CP-ABE encryption for vehicle information
        \item Shamir's Secret Sharing for emergency data (5,3 threshold)
        \item Keccak256 hashing for ride agreements
        \item Merkle trees for batch verification
    \end{itemize}
    
    \item \textbf{Layer 3 - Verification \& Governance}:
    \begin{itemize}
        \item Zero-Knowledge Proofs for driver verification
        \item DPoS delegate voting for disputes
        \item On-chain audit trail for immutable records
    \end{itemize}
\end{enumerate}

\subsection{Authentication Flow}

The authentication process follows these steps:

\begin{enumerate}
    \item User enters institutional email (@iiitkottayam.ac.in)
    \item System generates 6-digit OTP with 10-minute expiry
    \item OTP sent via email using SMTP (Gmail)
    \item User verifies OTP (max 3 attempts)
    \item Upon verification, user completes registration
    \item Password hashed using bcrypt (10 salt rounds)
    \item User connects MetaMask wallet
    \item System mints Soulbound Token to wallet address
\end{enumerate}

\subsection{Authorization Mechanism}

Access control is implemented through middleware:

\begin{lstlisting}[language=Java, caption=Authentication Middleware]
const authenticateToken = async (req, res, next) => {
    const userId = req.headers['x-user-id'];
    if (!userId) {
        return res.status(401).json({
            error: 'Authentication required'
        });
    }
    const user = await User.findByPk(userId);
    if (!user || user.isBanned) {
        return res.status(403).json({
            error: 'Access denied'
        });
    }
    // Verify SBT on-chain for sensitive operations
    const hasSBT = await userIdentityContract.verifySBT(user.walletAddress);
    req.user = user;
    req.hasSBT = hasSBT;
    next();
};
\end{lstlisting}

\subsection{Data Validation}

Input validation occurs at multiple levels:

\begin{itemize}
    \item \textbf{Client-side}: React form validation with real-time feedback
    \item \textbf{API Level}: Express middleware validation
    \item \textbf{Database Level}: Sequelize model constraints
    \item \textbf{Smart Contract Level}: Solidity require statements
\end{itemize}
